\section{Real-axis reconstruction and physical conformal currents}
\label{sec:reconstruction}

By Theorem~\ref{P:return}, the extracted scalar wave is the entire
trace polarization: $S=l$.  In this Section we write $L$ for that accessible
entry.  Until the joint domain is constructed below, this notation denotes
its vacuum wave and its analytic products, rather than an operator on
arbitrary states.  In particular, the vacuum-wave division by $M^2$ is
never extended to another leg.  The ingredients from Section~\ref{sec:extraction} are
bounded local approximations, analytic vector products of every finite
length, their reflected Hilbert splittings, and gluing at spacelike ties.
The leg equation is
\begin{equation}\label{R:leg-equation}
 \Box_i W(\ldots,L(x_i),\ldots)
       =W(\ldots,\Theta(x_i),\ldots).
\end{equation}
The remaining task is to establish joint tempered boundaries, their
ordered spectra, and affiliation with the original local algebras.

The construction has three stages.  First, we obtain real-axis
boundaries for pure scalar products and then for every mixed ordering.
Second, their Hilbert splittings provide a common invariant domain,
and bounded approximants give affiliation in the fixed net of
Assumption~\ref{ass:framework}.  Only after these steps do we regard
$L$ as a physical field and use its improvement to define the local
conformal Ward derivations.

The first stage proves the following joint boundary statement.  Its
scalar and vector conclusions are both needed: scalar positivity
defines insertion operators, while vector continuity identifies their
domain with actual vectors in the original Hilbert space.

\Needspace{6\baselineskip}
\begin{theorem}\label{R:mixed-products}
For every finite ordered list of $L$ and original physical fields,
the mixed analytic scalar product is the Fourier--Laplace continuation
of a tempered Wightman distribution, with the ordered spectrum,
adjunction, locality, and
similarity covariance, with dilations acting on the original field
labels.
The analytic vector products have tempered Hilbert-valued boundaries
with the vector spectrum.  Their reflected splittings hold for joint
tests and give all joint Wightman positivity inequalities in the
original Hilbert space.
\end{theorem}

We first prove this for homogeneous original slots.  Pure scalar
vectors give the initial mixed orderings; an adjacent-interchange
argument then reaches every ordering.  After constructing the vector
boundaries and their splittings, finite multilinearity restores the
arbitrary original field labels under the same joint test kernels.

\subsection{Uniformity conventions and pure scalar products}

All assertions below concern a fixed finite list of fields and compact
real coordinate sets.  For a test supported in a compact set $K$, put
\begin{equation}\label{R:seminorm}
 p_{K,m}(f)=\max_{|\alpha|\leq m}\sup_K|\partial^\alpha f|.
\end{equation}
A distribution estimate is uniform on the unit ball of one such seminorm;
the finite integer $m$ is allowed to increase when a prescribed accuracy
increases.  For boxes of diameter $\varepsilon$, the same convention
uses rescaled tests and hence introduces only powers of
$\varepsilon^{-1}$.  Imaginary directions range over compact subsets of
open ordered cones.  Constants can depend on those compact sets, the
field list, and the requested derivative order, but never on the dyadic
approximation parameter or on $0<\varepsilon\leq\varepsilon_0$.

For clarity we record the quantitative content of the bounded
approximation result.  For the product of a fixed number of approximants,
there are fixed $A,Q<\infty$ such that its holomorphic norm is at most
$C\varepsilon^{-A}R^Q$.  On a separated real placement patch its dyadic
difference satisfies, for every $k$,
\begin{equation}\label{R:approx-data}
 \|F_{2R}-F_R\|\leq C_k\varepsilon^{-B_k}R^{-k},
 \qquad
 \|F_R\|\leq C\varepsilon^{-A}R^Q.
\end{equation}
The first estimate follows by telescoping and moving the differing
bounded factor to the vacuum; the finite seminorms for it depend on $k$.
The exponent $Q$ in the second estimate does not.  Both estimates include
the adjoint waves.  These are precisely the shrinking-box estimates of
Proposition~\ref{E:bounded-approximation}; they impose no energy bound on
an unknown field operator.

The two exponents in \eqref{R:approx-data} play different roles.
Increasing the desired accuracy changes $k$ and the required test
seminorm, while the polynomial norm exponent $Q$ remains fixed.
This is what permits a summable dyadic estimate after a positive
fraction of the accuracy has been lost in analytic continuation.

We will obtain spectral support from scaling and local boundary
regularity.  The exact homogeneous case suffices for pure products.
Later we will also need the version in which a scaling error already
has the required spectral support.

\Needspace{6\baselineskip}
\begin{lemma}[Scaling and tube spectrum]\label{R:scaling-spectrum}
Let $C\subset\R^N$ be an open convex cone and set
\[
 K=C^*=\{\xi:\xi\cdot y\geq0\ \hbox{for every }y\in C\}.
\]
Use Fourier waves $e^{ix\cdot\xi}$.  Let $u$ be a tempered distribution
with values in a Hilbert space, locally near the origin the
distributional boundary of a holomorphic function on a wedge with
imaginary directions in $C$.  Suppose that, for a real $d$ and
unitaries $V(\lambda)$, $\lambda>0$, every scaling defect
\[
 \lambda^d u(\lambda x)-V(\lambda)u(x)
\]
has Fourier support in $K$.  Then $\supp\widehat u\subset K$.
In particular this holds for exact scaling covariance, with zero
defect.  For scalar distributions one may take $V(\lambda)=1$.
The wedge-boundary hypothesis may be replaced by rapid decay of
localized Fourier transforms near the origin in every direction
outside $K$.
\end{lemma}

\begin{proof}
A distributional wedge boundary has no analytic, hence no smooth,
wavefront at the origin in a direction outside $K$.  This follows
by shifting a localized Fourier test a small distance in a strict
direction of $C$ with negative frequency pairing.  The usual
distributional boundary bounds on smaller approach wedges justify
the shift.  The argument applies in Hilbert norm, or after
scalarization uniformly on the unit ball.

Fix a smooth Fourier test supported in a compact subset of
$\R^N\setminus K$.  Its pairing with $\widehat u$ after dilation
to frequency $\Lambda$ decreases faster than every power of
$\Lambda$.  To pass from the local wavefront estimate to this
unlocalized pairing, insert a cutoff equal to one near the origin.
The local part has the asserted rapid decay.  Outside the cutoff,
the rescaled Schwartz inverse Fourier test and all its derivatives
decrease rapidly in every Schwartz seminorm, so temperedness
controls the remainder.  A finite covering treats all directions
in the fixed frequency support.

The scaling defect has zero pairing with that test and with every
one of its positive dilates, since $K$ is a cone.  The remaining
scaling identity makes the norm of the dilated pairing a fixed
power of $\Lambda$ times the norm of the original pairing;
$V(\lambda)$ does not change that norm.  Rapid decay forces the
original pairing to vanish.  Every compact Fourier test outside
$K$ is allowed, proving the assertion.
\end{proof}

\Needspace{6\baselineskip}
\begin{proposition}\label{R:pure-products}
Every pure $L$ analytic vector product has a tempered Hilbert-valued
boundary with the ordinary ordered vector spectrum.  Its Fourier--Laplace
continuation is the accessible analytic product.
\end{proposition}

\begin{proof}
First, the product has scalar-primary covariance under sufficiently
small simultaneous special transformations.  The geometry in
Lemma~\ref{P:common-wedge} applies to every fixed finite family of
source and destination patches: its proof uses only finitely many
strict wedge inequalities.  Choose the conjugate null axis $m$ separated from the
finitely many flow directions, take $m\cdot x$ sufficiently negative,
and then add a sufficiently large positive multiple of $m$ to make
all the required opposite inequalities strict.  A small double cone
$O_-$ about that point is spacelike to every containing wedge and to
the destination patches.  Local vacuum cyclicity makes
$\{C\Omega:C\in\cM(O_-)\}$ total.

For a fixed finite product length, reparametrize the source tests by
a small null flow and its inverse primary factors.  Write $F_g^{(n)}$
for the limit of the resulting simultaneously conjugated bounded
products and $F^{(n)}$ for the accessible destination product.  The
former limit exists because the source product converges and the
common conjugating unitary is bounded.  For $C\in\cM(O_-)$, move
its first bounded factor past $C$ and take the vacuum and
adjoint-vacuum limits.  The canonical one-wave action and induction
on the length give
\begin{align}\label{R:all-length-recursion}
 \langle C\Omega,F_g^{(n)}\rangle
   &=\langle L(\bar f_1)\Omega,C^*F_g^{(n-1)}\rangle\\
   &=\langle L(\bar f_1)\Omega,C^*F^{(n-1)}\rangle
    =\langle C\Omega,F^{(n)}\rangle.\notag
\end{align}
Every transformed one-vacuum wave is already the accessible $L$
wave, so this induction uses no finite-length interpolation.  Totality
proves the covariance identity on the real spacelike patch, and
boundary uniqueness proves it in the tube.  Compositions of small
null flows give small special transformations.  For any fixed real
compact set and strict approach subcone, choose the transformations
small both there and on a finite interior path to the initial patch.
Their unitary factors have norm one.  Thus the identity is uniform
near that compact set; localization under arbitrary large conformal
transformations is not required.

Thus covariance has been obtained separately for every finite pure
product length.  The next step uses it to transport the known
boundary estimates for products of $\Theta$ into enough differential
estimates to recover the original pure $L$ product.

Here is an algebraic device that converts this covariance into a
boundary estimate.  On the ambient null cone use
\[
 X=(X^+,X^-,X^\mu),\qquad q(X)=X^\mu X_\mu-X^+X^-,
 \qquad X=\lambda(1,x^2,x).
\]
Here $F$ is the Hilbert-valued analytic vector product.  Its lift is
Hilbert-valued and homogeneous of degree $w=-2$ in each slot;
the differential identities below may be tested against arbitrary
Hilbert vectors, with estimates uniform on their unit ball.  On degree
$w$ homogeneous functions, define the flat Thomas derivative
\begin{equation}\label{R:Thomas}
 \mathscr D_A(w)=(2+2w)\partial_A-X_A\partial^2_{\mathrm{amb}}.
\end{equation}
Successive derivatives use the new degree in their slot.  The product
rule, $\partial^2_{\mathrm{amb}}q=12$, and Euler's identity show that
these derivatives preserve the quadratic ideal, commute successively,
and have zero successive trace.  They therefore act intrinsically on
the cone.  The same calculation gives
\begin{equation}\label{R:Thomas-contraction}
 X^A\mathscr D_A(w)=w(2+2w)\quad\pmod{q}.
\end{equation}
These are the flat ambient identities underlying the Thomas operator;
see also \cite{GoverWaldron2009}.  None of the factors in
\eqref{R:Thomas-contraction} vanishes at $w=-2,-3,\ldots$.

In particular, contraction can undo each successive derivative at
the degrees encountered here.  The weight in
\eqref{R:Thomas} must be lowered after each differentiation in that
slot; using a fixed weight throughout would not give these intrinsic
commuting operators.

Take the null direction $a_\infty$ representing chart infinity, and
extend the lifted function independently of $X^-$.  Then
$a_\infty\cdot\mathscr D$ is a nonzero smooth chart factor times
$\Box_x$.  Consequently
\begin{equation}\label{R:all-slot-derivative}
 G_aF:=\prod_{i=1}^n(a\cdot\mathscr D_i)F
\end{equation}
has polynomial boundary growth for $a=a_\infty$, by the ordinary
all-$\Theta$ product and \eqref{R:leg-equation}.  Primary transport
proves the same statement for $a$ in an open real null-cone patch about
$a_\infty$.  Finite derivatives preserve these bounds after shrinking
the approach subcone, by Cauchy's inequalities.

For this transport, write $\widetilde F$ for the homogeneous lift.
For each fixed sufficiently small transformation $g$, its local
covariance is $\widetilde F(gX)=U_g\widetilde F(X)$, with the
primary chart factors absorbed in the lift.  The flat ambient
chain rule gives $G_{ga}\widetilde F(gX)=U_gG_a\widetilde F(X)$.
The unitary $U_g$ is fixed as the insertion coordinates vary, so
the coordinate derivatives act only on $\widetilde F$.  Taking
$a=a_\infty$ transfers the known bound to the null direction
$ga_\infty$.  The shifted homogeneous weights supply smooth
chart factors bounded on the chosen compact set.  Scalarization
replaces a testing vector by its unitary image and therefore keeps
the Hilbert norm estimate uniform.

We give the finite algebraic argument needed to recover $F$.  In
constant complex variables $\xi_{iA}$, let $I$ be the ideal generated by
$\xi_i^2$ and by $\prod_i(a\cdot\xi_i)$ for all $a$ in that patch.
If a common zero had every $\xi_i\ne0$, the patch would be covered by
the finitely many proper hyperplane sections $a\cdot\xi_i=0$.
The real null patch is Zariski dense in the irreducible complex null
cone and spans the ambient space, so this is impossible.  Thus every
common zero has $\xi_i=0$ for some $i$.

For each of the finitely many monomials
$m=\prod_i\xi_{iA_i}$, the Nullstellensatz gives $m^{e_m}\in I$.
Let $e$ dominate their exponents.  Every monomial of degree at least
$6(e-1)+1$ in each slot is divisible by one of the $m^e$ and belongs
to $I$.  Noetherianity gives finite generator representations, and
separate homogeneity in each slot gives multihomogeneous coefficient
representations.  Substitute the successively commuting operators
\eqref{R:Thomas}.  Trace terms vanish; the remaining terms are finite
derivatives of \eqref{R:all-slot-derivative}.  Hence sufficiently high
derivative tensors of $F$ have polynomial growth.  Contracting their
indices with the $X_i$ and using \eqref{R:Thomas-contraction} repeatedly
recovers $F$ with a nonzero constant factor.  The chart factors are
bounded on the compact set.  This proves moderate growth of $F$.

The algebraic recovery used only a finite number of null directions
and derivatives for this fixed product length.  Its constants may
depend on that length; no estimate uniform over all lengths is needed
for constructing the polynomial field domain.

The moderate-growth tube boundary theorem, applied weakly and then
uniformly in unit Hilbert vectors, gives a Hilbert-valued distribution
locally.  Its exact similarity covariance globalizes this boundary
and makes it tempered: rescaling dyadic annuli to a fixed compact set
changes a finite distribution seminorm by only a power of the radius,
and the dilation unitary leaves its norm unchanged.  Lemma~\ref{R:scaling-spectrum}
gives its ordered vector spectrum.  Fourier--Laplace uniqueness
then identifies its continuation with the original product.
\end{proof}

\subsection{Boundary values and a distribution criterion}

We specify the analytic facts used in the mixed-product argument so
that no growth assertion is concealed in a boundary-value notation.
The point of the distinction is that holomorphy alone will initially
give boundary hyperfunctions for the coefficients of a causal wave
ambiguity.  We will prove a quantitative Gaussian estimate before
using any of those coefficients as an ordinary distribution.
The boundary $b f$ of a holomorphic function on a local convex wedge
exists as a hyperfunction without a growth assumption.  Restriction
to a smaller approach cone is compatible with this boundary map; the
map is injective for the local germ on a connected single wedge and
agrees with ordinary distributional boundary values whenever those
exist.  We fix the Euclidean coordinate volume form when identifying
scalar hyperfunctions with the corresponding top-form convention.  Compactly
supported hyperfunctions are analytic functionals carried by their
support.  Thus, for each complex neighborhood $K_\gamma$ of the
carrier, a representative $u$ satisfies
\begin{equation}\label{R:carrier}
 |\langle u,h\rangle|\leq C_\gamma\sup_{K_\gamma}|h|,
 \qquad h\text{ holomorphic near }K_\gamma.
\end{equation}
Komatsu constructs the injective boundary map on a Stein complex
neighborhood with an open convex imaginary cone
\cite[Section~1]{Komatsu1973}; equations (1.1) and (1.5) on
pp.~253--254 give injectivity and compatibility with smaller cones.
Small complex boxes are Stein, so restriction to a smaller real patch
gives the local germ statement used here.  The comparison on p.~257
identifies polynomial-growth boundaries with distributional ones.
Martineau's Theorems~1--2 give the analytic-functional realization
\cite{Martineau1961}; the carrier estimate is in Section~1(a)--(b)
and Proposition~1 of that paper.  These constructions alone give no
polynomial growth estimate.  In the analytic-functional notation
used below, the canonical distribution embedding is described in
\cite[Section~4, equation~(4.3)]{HoepfnerRagognette2021}.

Use real Euclidean coordinates for the Gaussian, with their
complex-bilinear extension, and the Fourier wave convention
$e^{ix\cdot\xi}$.  For a compact analytic-functional representative put
\begin{equation}\label{R:FBI}
 \mathcal T_\Lambda u(w,\eta)
   =\left\langle u(z),
       e^{i\Lambda(w-z)\cdot\eta-\Lambda(w-z)^2/2}\right\rangle.
\end{equation}
Only real centers $w$ and real frequencies $\eta$ are used for the
estimates below.

\Needspace{6\baselineskip}
\begin{lemma}\label{R:FBI-criterion}
Suppose that, near a real point $w_0$, a hyperfunction $u$ has a compact
analytic-functional representative such that
\begin{equation}\label{R:FBI-polynomial}
 |\mathcal T_\Lambda u(w,\eta)|\leq C(1+\Lambda)^N
 \quad(w\in U, |\eta|=r_*,\ \Lambda\geq1)
\end{equation}
for some neighborhood $U$ of $w_0$ and some $r_*>0$.
Then $u$ is a distribution near $w_0$.

If $u=b f$ from a wedge, its transform near $w_0$ can be computed by
integrating $f$ on a sufficiently small fixed interior translate of
a local real box, with an exponentially small error.  The box and
height can be chosen uniformly for $w$ in a smaller box and $\eta$
in a fixed compact set.
\end{lemma}

\begin{proof}
The first assertion is the distribution case of the generalized FBI
criterion of Hoepfner--Ragognette
\cite[Theorem~6.3 and Section~8.1]{HoepfnerRagognette2021}.
In their notation take the positive quadratic polynomial
$p(v)=|v|^2/(2r_*)$ and $\xi=\Lambda\eta$, $|\eta|=r_*$.
Then $|\xi|p(w-z)=\Lambda(w-z)^2/2$, so their transform is
\eqref{R:FBI} up to an irrelevant fixed normalization.  Their local
inverse decomposition separates a distribution, obtained from the
polynomially bounded part, from an analytic remainder.  Finite
covering of the frequency sphere makes the polynomial estimate
uniform.  The compact representative is immaterial because changing
it away from $w_0$ changes this estimate by an exponentially small
term.

For the second assertion, use the fixed-height integration
representation of a wedge boundary value
\cite[Remark~4.1]{HoepfnerRagognette2021}, using a rounded box with
smooth boundary.  The difference is an analytic
functional carried in an arbitrarily thin complex neighborhood of
the boundary of the chosen real box.  This is also obtained from the
product Cauchy representation: entries that extend across one
coordinate cancel under contour displacement.  Axes strictly inside
the wedge give the requisite product half-domain description.
Let the distance between the smaller box of centers and that boundary
be $a>0$.  For $z=x+iy$, $|\eta|\leq M$, the real exponent in the
Gaussian is
\begin{equation}\label{R:Gaussian-exponent}
 \Lambda\left(y\cdot\eta+\frac{|y|^2}{2}
                    -\frac{|w-x|^2}{2}\right).
\end{equation}
Choose all side contours with $|y|\leq t$ and then the carrier
neighborhood with width $\gamma$ so small that
$(t+\gamma)M+(t+\gamma)^2/2<a^2/4$.
Equation~\eqref{R:carrier} bounds their contributions by $C e^{-c\Lambda}$.
This is the localization estimate of
\cite[Proposition~4.5]{HoepfnerRagognette2021}.  Heights and carrier
neighborhoods are fixed before $\Lambda\to\infty$; no integrability
of $f$ down to the real edge has been assumed.
\end{proof}

A related fact for a tempered distribution will be useful.  Rapid
Gaussian decay near the origin in a frequency cone implies rapid
decay when its Fourier transform is tested at large frequencies
strictly inside that cone.  One can use the smooth wavefront
characterization, or directly integrate the Gaussian inversion
against that frequency test.  Away from the origin, integrations by
parts in the frequency variable gain arbitrary powers; a derivative
of the Gaussian costs at most a square-root power beyond the fixed
distribution order.  This gives the alternative Fourier-decay
hypothesis in Lemma~\ref{R:scaling-spectrum}.

\subsection{A quantitative boundary estimate at one merged face}

Until the finite-multilinearity passage below, every original-field
slot in the mixed-product argument is homogeneous, including every
spectator slot.  Original derivative slots carry their shifted
dimensions; derivatives of $L$ will follow by differentiating its
reconstructed products.
We fix an arbitrary finite list and an arbitrary ordering; neither its
length nor the set of lists is bounded.  The constants and distribution
orders may depend on this list.  We shall first complete both scalar
and Hilbert-valued reconstruction for these homogeneous lists.

The scalar induction is simultaneous over every finite homogeneous
old-field list and is indexed only by the number of $L$ entries.
Its zero stage is the original Wightman family.  At the next stage,
orderings with every $L$ on the right are tempered: the analytic
splitting \eqref{E:accessible-splitting} pairs an ordinary
homogeneous old bra with the pure vector boundary of
Proposition~\ref{R:pure-products}.  This continuous bilinear pairing
on the conjugate bra Hilbert space and the ket Hilbert space gives a
joint Schwartz distribution by the kernel theorem.  The reversed
bra vector cone, the ket vector cone, and translation invariance
supply its scalar ordered cones.

It remains to move $L$ entries through old entries.  In each adjacent
interchange, applying $\Box$ in the $L$ slot replaces it by the
homogeneous old field $\Theta$.  Both source orderings therefore
belong to the preceding induction stage, however many old spectators
are present.  We prove the interchange step below before iterating it.

Let $W_1$ be a mixed scalar ordering already known to be tempered,
and obtain $W_2$ by interchanging an adjacent $L(x)$ and homogeneous
original physical field
$A(y)$.  Write $r=x-y$.  The variable $z$ consists of $y$ and all
spectator coordinates in their merged order.  Holding $z$ in a compact
subset of its strict tube keeps every gap except the new relative
gap bounded away from the boundary.  Replacing this $L$ by $\Theta$
gives two already known tempered sources, with the homogeneous
dimension-four field $\Theta$ in that slot,
$S_i=\Box_rW_i$.

In these coordinates the momentum at the merged $y$ slot is the sum
of the two adjacent momenta; all other spectator momenta are unchanged.
Write $C_z$ for the strict merged approach cone and $K_z=C_z^*$ for
its closed spectral cone, in the Fourier convention of
Lemma~\ref{R:scaling-spectrum}.  A strict relative direction $v_r$
of the target order, sufficiently small compared with a fixed merged
direction $v_z$, makes $(v_r,v_z)$ a strict full target direction.
These are different cones: operations in $r$ preserve $K_z$, while
the missing relative spectral condition will require a further argument.

Here the relative gap tends to the boundary while the merged gaps
remain strict.  The estimate must therefore keep a positive analytic
propagation exponent as the relative gap shrinks.  Once this face
estimate is available, the wave equation will carry it from a thin
real-time strip to every compact relative-coordinate set.

\Needspace{6\baselineskip}
\begin{lemma}\label{R:face-bound}
For compact real $r$-sets, compact strict merged $z$-sets, and compact
sets of relative imaginary directions $v$ of the target order, there
are $C,N<\infty$ such that
\begin{equation}\label{R:face-power}
 |W_2(r+i\varepsilon v;z)|\leq C\varepsilon^{-N},
 \qquad 0<\varepsilon\leq\varepsilon_0.
\end{equation}
The same statement holds for every fixed finite derivative order.
Consequently $W_2$ has an ordinary distributional $r$-boundary
$F_0(r;z)$, holomorphic as a distribution-valued function on the strict
merged tube.
\end{lemma}

\begin{proof}
We first work with $|r^0|\leq C_0\varepsilon$ and bounded
$\mathbf r$.  Shift $\mathbf r$ outward by at most $C_1\varepsilon$
so that the resulting $r_*$ has spacelike clearance at least
$4c\varepsilon$.  Put the real merged placements $z_*$ sufficiently
far apart on a common time slice.  Their spacelike margins, including
separation from both members of the pair, can be chosen of order one.
The tests for the two nearby slots have diameter $c\varepsilon$;
all tests may be taken of that size.  The bounds
\eqref{R:approx-data} hold on the resulting separated patch.

Here $F_R$ denotes the scalar vacuum expectation of the bounded
product \eqref{E:bounded-product}.  On the vector tube, vacuum
invariance removes its initial factor $U(z_1)$.  The resulting
expression depends only on consecutive gaps and extends holomorphically,
with the same polynomial norm bound, to the scalar tube where those
gaps are strict.  Thus this scalar face estimate does not require the
first-gap condition that remains necessary for vector products.

Write $\eta=\operatorname{Im}z$.  Consider the complex line
\begin{equation}\label{R:anisotropic-line}
 \zeta\longmapsto
     (r_*+\varepsilon\zeta v,\ z_*+\zeta\eta),
       \qquad \operatorname{Im}\zeta>0.
\end{equation}
For sufficiently small $\varepsilon$ its direction is in the full
target cone.  There is a real interval $[-a,a]$ on which the
approximant supports remain spacelike, with $a>0$ independent of
$\varepsilon$: the relative displacement is of size $\varepsilon$
and the corresponding clearance is also of that size, whereas all
merged displacements and margins are of order one.  The harmonic
measure of this interval at $\zeta=i$ is
$\alpha=2\arctan(a)/\pi>0$.  The two-constants theorem applied to
$\log|F_{2R}-F_R|$ gives
\[
 |F_{2R}-F_R|
 \leq C_k\varepsilon^{-\alpha B_k-(1-\alpha)A}
           R^{-\alpha k+(1-\alpha)Q}.
\]
Slightly reducing the clearances gives this estimate on a complex
$r$-ball of radius $c_1\varepsilon$ and a $z$-ball of fixed radius.

Returning the spatial shift takes a bounded number of
$\varepsilon$-balls.  At fixed complex $r$, moving the merged
variables to any prescribed merged compact takes a bounded number
of fixed-radius $z$-balls.  Moving $y$ moves $x=y+r$ with it, so the
small internal gap does not restrict those $z$-radii.  More precisely,
take the compact convex hull, inside the strict merged tube, of the
initial placements and the prescribed compact set.  Its gaps have a fixed positive
margin; the two exterior gaps adjacent to the pair keep that margin
when the internal displacement is $O(\varepsilon)$.  A finite cover
therefore supplies radii and a number of balls independent of
$\varepsilon$.  Take overlapping
balls with radii $\rho,2\rho,4\rho$ inside the tube.  Each three-circle
step retains exponent $1/2$.  Thus a fixed $\sigma>0$, independent of
$\varepsilon$, gives
\begin{equation}\label{R:dyadic-face}
 |F_{2R}-F_R|
   \leq C_k\varepsilon^{-E_k}
             R^{-\sigma k+(1-\sigma)Q},
 \qquad E_k=\sigma\max(B_k,A)+(1-\sigma)A.
\end{equation}
Choose one finite $k$ with $\sigma k-(1-\sigma)Q>1$ and then fix its
finite test seminorm.  Summing dyadic differences gives a fixed
inverse power of $\varepsilon$.  If localization starts at a dyadic
$R_0(\varepsilon)$ bounded by a power of $\varepsilon^{-1}$, its
initial norm contributes another finite power.  The construction of
Proposition~\ref{E:bounded-approximation} permits this choice.  The
accuracy order is never increased while $\varepsilon$ tends to zero.

This order of choices is essential: the continuation geometry fixes
$\sigma$, then one chooses $k$, and only then lets the relative gap
shrink.  The resulting $E_k$ is a fixed finite exponent in the face
bound.

For completeness, unsmearing this estimate does not involve point
evaluation of a bounded approximant.  The limit already has a
holomorphic kernel by Proposition~\ref{E:accessible-products}.  Choose
a smooth radial bump $\rho$ on $\C^4$ of integral one, supported in
a sufficiently small ball.  Holomorphic mean value in each of $p$
slots gives, writing $Z$ for the full slot coordinates,
\begin{align}\label{R:unsmear}
 W_2(Z)
 &=\int_{(\C^4)^p}\prod_i\varepsilon^{-8}
               \rho(w_i/\varepsilon)W_2(Z+w)\,\dd w\\
 &=\int \dd b\;
     W_2(f_{\varepsilon,b_1},\ldots,f_{\varepsilon,b_p};
                     Z+i\varepsilon b).\notag
\end{align}
Here $f_{\varepsilon,b}(u)=\varepsilon^{-4}\rho(u/\varepsilon+ib)$,
and $b$ ranges over a fixed compact set and
$p_{K,m}(f_{\varepsilon,b})\leq C_m\varepsilon^{-4-m}$.
The finite seminorm estimate just proved therefore bounds the last
integral by a power.  Cauchy's inequalities in smaller
$\varepsilon$-balls also bound all fixed derivatives.  This proves
\eqref{R:face-power} initially in the thin real-time strip.

Finally fix $z$ and the imaginary displacement $i\varepsilon v$.
The function
$u_\varepsilon(t,\mathbf r)=W_2((t,\mathbf r)+i\varepsilon v;z)$
satisfies the ordinary inhomogeneous wave equation in the real
variables.  Its initial value, spatial derivatives, and first time
derivative have the bounds just obtained.  Its source is the
Fourier--Laplace continuation of the inductively tempered $S_2$,
so it has a power bound of its own at this relative gap.  Kirchhoff's
formula on $|t|\leq T$, $|\mathbf r|\leq B$ gives
\begin{align}\label{R:wave-estimate}
 \|u_\varepsilon\|_{T,B}
 &\leq \|u_\varepsilon(0)\|_{B+T}
   +T\|\nabla u_\varepsilon(0)\|_{B+T}
   +T\|\partial_tu_\varepsilon(0)\|_{B+T}
   +\tfrac12T^2\|S_2\|_{T,B+T},
\end{align}
with an immaterial fixed norm constant.  This proves the bound on
every real compact set.  Derivatives obey the same estimate.
Moderate-growth boundary values now give $F_0$; the estimates uniform
on $z$-compacts give holomorphy and locally uniform finite distribution
order in $r$.
\end{proof}

\subsection{Causal ambiguity and recovery on the wrong massless branch}

The source difference $S_2-S_1$ vanishes for real spacelike $r$ by
inductive locality.  Split it into distributions supported in
$J^+(0)$ and $J^-(0)$ by a partition away from $r=0$ and finite Taylor
subtraction at that point.  Subtraction acts only in $r$, so all
merged $z$-spectra are preserved.  Convolve these two parts with the
retarded and advanced fundamental solutions, respectively.  The
cone-supported convolutions are proper on bounded sets and define
tempered distributions; their Schwartz seminorms involve only
finitely many seminorms of the sources.  Adding their sum to $W_1$
gives a tempered distribution $T_0$ satisfying
\begin{equation}\label{R:causal-solution}
 \Box_r T_0=S_2,\qquad
 \supp_r(T_0-W_1)\subset J(0),
\end{equation}
with the merged spectra.  Let the same symbol denote its
Fourier--Laplace continuation in $z$, valued in $r$-distributions.

At fixed strict $z$, the difference $H_*=F_0-T_0$ solves the
homogeneous wave equation and is supported in $J(0)$; spacelike tie
gluing proves the support assertion.  A homogeneous wave distribution
has distributional Cauchy data.  At time zero those data are supported
at the spatial origin, so each is a finite sum of derivatives of a
delta function.  Cauchy uniqueness consequently gives
\begin{equation}\label{R:wave-ambiguity}
 H_*(r;z)=\sum_{j=1}^m C_j(z)P_j(\partial_r)\Delta_0(r).
\end{equation}
The $P_j$ may be chosen homogeneous and linearly independent modulo
the wave quadratic, with $\Delta_0$ the massless commutator function.
One fixed finite list works on the connected merged tube.  Indeed,
on a relatively compact open subset the distribution order is
uniform, so sufficiently high spatial monomials annihilate both
Cauchy data.  These annihilators continue holomorphically to the
whole tube.  Cauchy restriction is continuous on homogeneous wave
solutions, and fixed dual tests recover the finitely many coefficients;
hence the $C_j$ are holomorphic.

The unknown part of the boundary has now been reduced to finitely
many scalar coefficient functions on the merged tube.  The next
lemma proves their distribution regularity by testing on the
massless branch excluded by the target ordering.  The causal
solution $T_0$ is already tempered and will supply the polynomial
bound in that test.

\Needspace{6\baselineskip}
\begin{lemma}\label{R:coefficient-distributions}
Every coefficient $C_j$ in \eqref{R:wave-ambiguity} has a local ordinary
distribution boundary in the merged variables.
\end{lemma}

\begin{proof}
Fix a real merged center.  Initially $bC_j$ is only a hyperfunction.
Choose a small real box about the center, a smaller box of centers
$w$, and a strict merged direction $v_z$.  Lemma~\ref{R:FBI-criterion}
represents the Gaussian testing by a fixed translate $it_0v_z$,
up to exponentially small side errors.

Choose a strict relative direction $v_r$ for the target ordering,
small enough compared with the merged gaps that $(v_r,v_z)$ is a
strict full target direction.  Choose finitely many nonzero null
covectors $k_a$ on the wrong branch, so $k_a\cdot v_r<0$.
The sample points range over a radially open patch of that branch;
they are not constrained to a fixed energy shell.
The radius variation will be used when separating coefficient
polynomials of different degrees; choosing directions only on one
fixed shell would not give that independence argument.
Then choose
$r_*>0$ sufficiently small that, uniformly for all samples and
$|\eta|=r_*$,
\begin{equation}\label{R:negative-phase}
 k_a\cdot v_r+\eta\cdot v_z\leq-c_0<0.
\end{equation}
Test \eqref{R:wave-ambiguity} in $r$ with
$\chi(r)e^{-i\Lambda r\cdot k_a-\Lambda r^2/2}$, where $\chi=1$
near zero, and in $z$ with the Gaussian
\eqref{R:FBI} on the fixed translated box.

For the $F_0$ term, deform the relative contour from zero to
$it_0v_r$, leaving $z$ at height $it_0v_z$.
Lemma~\ref{R:face-bound} justifies the initial endpoint.  On the final
central contour the real exponent, apart from negative spatial
Gaussians, is at most
\[
 \Lambda\left[-t_0c_0+
                  \tfrac12t_0^2(|v_r|^2+|v_z|^2)\right].
\]
Choose $t_0$ so small that this is $\leq-c\Lambda$.  The continued
$W_2$ is bounded on that fixed interior compact, so this term is
exponentially small.

We specify the cutoff error, since the relative contour starts at a
distributional face.  Let $N$ be an inverse-height growth exponent
from \eqref{R:face-power}, increased to cover the derivatives used
below.  Replace $\chi(x)$ on the moving contour by its almost-analytic
Taylor lift
\[
 \chi^{[M]}(x+iy)=
     \sum_{|\alpha|\leq M}(iy)^\alpha\partial^\alpha\chi(x)/\alpha!,
 \qquad M>N+1.
\]
Its antiholomorphic derivative is $O(|y|^M)$, supported in the
transition region.  There $|x|$ is separated from zero.  By decreasing
$t_0$ before taking the large-parameter limit,
\eqref{R:Gaussian-exponent} is $\leq-c'\Lambda$ throughout the
transition, even after all linear phase costs.  Stokes errors are
bounded by a polynomial in $\Lambda$ times
$e^{-c'\Lambda}\int_0^{t_0}t^{M-N}\,\dd t$.
The same construction applies to any smooth central cutoff in $z$.
For the merely hyperfunctional coefficient boundary we instead use
the fixed-height analytic-functional remainder of
Lemma~\ref{R:FBI-criterion}; we do not multiply that boundary by a
smooth cutoff.

The $T_0$ term is polynomially bounded in $\Lambda$: move its merged
contour back to the real boundary and use its tempered seminorm.
Derivatives of the Gaussian and cutoffs cost only powers of
$\Lambda$.  Its contour errors obey the same estimate, now with the
known finite distribution order of $T_0$.

It remains to check that these tests recover every coefficient.  Put
$d_j=\deg P_j$ and
\[
 M_{aj}(\Lambda)=
  \left\langle P_j(\partial)\Delta_0,
       \chi(r)e^{-i\Lambda r\cdot k_a-\Lambda r^2/2}\right\rangle.
\]
The Fourier transform of $\Delta_0$ is a nonzero normalization times
$\operatorname{sgn}(p^0)\delta(p^2)$.  Gaussian Fourier integration,
with $p=\Lambda k_a+\sqrt\Lambda q$, gives
\begin{equation}\label{R:shell-asymptotic}
 \Lambda^{3/2-d_j}M_{aj}(\Lambda)
       \longrightarrow c(k_a)P_j(ik_a),\qquad c(k_a)\ne0.
\end{equation}
Indeed $\dd p$ contributes $\Lambda^2$, the delta function
$\Lambda^{-3/2}$, and the four-dimensional Gaussian Fourier factor
$\Lambda^{-2}$.  The remaining integral is a nonzero Gaussian
integral on $2k_a\cdot q=0$, with the fixed sign of the branch.
The cutoff error is exponentially small by the temperedness of
$P_j(\partial)\Delta_0$.  The nonzero factor $c(k_a)$ can depend on
the row.

The evaluation matrix $P_j(ik_a)$ has full column rank for a suitable
finite sample set.  Otherwise a nonzero linear combination would
vanish on an open real patch of one null branch, hence on the
irreducible complex quadratic cone, contradicting independence
modulo the wave quadratic.  Radial openness matters for polynomials
of different degrees.  Selecting an invertible square submatrix,
\eqref{R:shell-asymptotic} bounds its inverse by powers of $\Lambda$.
The tested identity therefore gives
\[
 |\mathcal T_\Lambda(bC_j)(w,\eta)|\leq C(1+\Lambda)^N
      \qquad(|\eta|=r_*),
\]
uniformly on the smaller box.  Lemma~\ref{R:FBI-criterion} proves
that every $bC_j$ is an ordinary distribution there.
This conclusion is local in the merged real variables.  Growth at
infinity and the merged spectral support still require the scaling
argument below.
\end{proof}

\subsection{Global growth, spectra, and completion of mixed products}

Local distribution boundaries do not by themselves imply global
temperateness.  Here exact similarities provide the missing control.
The comparison solution $T_0$ need not share the exact homogeneity
of $W_2$.  Its scaling defect will instead appear as an explicitly
tempered error, which is enough both for dyadic growth estimates and
for excluding frequencies outside the merged spectral cone.
If $W_2$ has $n_L$ entries $L$ and homogeneous original entries of
dimensions $\Delta_1,\ldots,\Delta_k$, put
\[
 d_W=2n_L+\sum_{i=1}^k\Delta_i.
\]
Here the old dimensions already include any derivative shifts.
The accessible analytic product has exact similarity covariance for
this homogeneous list, by covariance of the one-wave data and uniqueness
of the multilinear accessible products.  Its distributional
$r$-boundary $F_0$ has the same exact scaling.  Put
$\delta_j=2+\deg P_j$, so
$P_j(\partial)\Delta_0(\lambda r)=\lambda^{-\delta_j}
P_j(\partial)\Delta_0(r)$.
Choose fixed compact $r$-tests $\varphi_j$ dual to the finite list
$P_j(\partial)\Delta_0$.  Exact scaling of $F_0$ and
\eqref{R:wave-ambiguity} show that
\begin{equation}\label{R:scaling-defect}
 \lambda^{d_W-\delta_j}C_j(\lambda z)-C_j(z)
 =\left\langle
      T_0(r;z)-\lambda^{d_W}T_0(\lambda r;\lambda z),
                   \varphi_j(r)\right\rangle_r.
\end{equation}
The right side has a tempered boundary with merged spectral support.
For $\lambda\geq1$ its Schwartz seminorm constants grow at most
polynomially in $\lambda$: scaling and the fixed compact extraction
tests act continuously on Schwartz space with such bounds.

Cover the complement of a compact $z$-ball by dyadic annuli and
rescale each annulus to a fixed compact set.  There the local
coefficient distribution has a fixed finite order.  Formula
\eqref{R:scaling-defect} bounds its value on every outer annulus by
a polynomial in the annulus radius and finitely many Schwartz
seminorms.  Summing against a Schwartz test proves that each $C_j$
is tempered.  This argument needs no exact homogeneity of $T_0$.

The scaling defect in \eqref{R:scaling-defect} has Fourier support
in $K_z$.  Apply Lemma~\ref{R:scaling-spectrum} to the now tempered
coefficient boundary, with $d=d_W-\delta_j$ and $V(\lambda)=1$.
Its local wedge boundary is supplied by
Lemma~\ref{R:coefficient-distributions}.  We obtain
$\supp\widehat C_j\subset K_z$.
Fourier--Laplace continuation and injectivity identify the original
$C_j$ with the corresponding tempered continuation; in particular
it now has ordinary polynomial bounds at ordered approaches.

It follows that $F_0$ is the merged continuation of the globally
tempered distribution
\begin{equation}\label{R:D0}
 D_0(r,z)=T_0(r,z)+\sum_j C_j(z)P_j(\partial_r)\Delta_0(r).
\end{equation}
It is exactly homogeneous, because its merged continuation $F_0$
is.  We next establish the full target spectrum, rather than just
the merged one.  For a covector strictly outside the full target
cone, choose a strict target imaginary direction $(v_r,v_z)$ with
negative pairing.  Local Gaussian testing near $(r,z)=0$ first
moves $z$ to a small height $it_0v_z$, using the already proved
merged temperateness.  At that face, $D_0=F_0$; move $r$ to
$it_0v_r$ using Lemma~\ref{R:face-bound}.  Intermediate fractions of
the latter move remain in the target tube.  The final contour is an
interior $W_2$ contour and has exponential Gaussian decay.  The side
errors are those already estimated in the proof of
Lemma~\ref{R:coefficient-distributions}.  This supplies the local
off-cone Fourier decay used in Lemma~\ref{R:scaling-spectrum}; exact
homogeneity of $D_0$ therefore gives its full target spectral support.
Its full Fourier--Laplace continuation agrees with $W_2$, since they
have the same $r$-boundary at fixed strict merged damping.

One adjacent interchange is therefore complete: it preserves
temperateness, recovers the full ordered spectrum, and retains
spacelike locality.  Iterate this step in the simultaneous
induction stated above.  At fixed number of $L$ entries, begin with
all of them on the right and move them to the prescribed positions.
Every source has one fewer $L$ and one additional old entry $\Theta$,
so it belongs to the preceding induction stage.  Repeat this for
every finite ordering of the old entries.  Tie gluing, followed by the
ordinary merged spectral boundary with a spacelike relative test,
then proves locality also for adjacent old-field interchanges, with
the prescribed grading.  Adjunction and exact similarity covariance
for each homogeneous list follow from those identities of the analytic
accessible products and uniqueness of their boundaries.

We next complete the Hilbert-valued assertion for the same homogeneous
lists.  The squared norm of a mixed analytic vector product is its
reflected scalar split.  A strict vector approach gives a strict
scalar split approach: the central gap is twice the first vector
damping, and the other gaps are the reflected vector gaps.  The scalar
polynomial bound therefore gives a polynomial norm bound for the vector
product.  The Hilbert-valued moderate-growth theorem supplies its local
boundary.  Its exact similarity covariance for this homogeneous list,
with a unitary dilation factor on the output, globalizes the boundary
by dyadic annuli and gives the vector support by
Lemma~\ref{R:scaling-spectrum}.  Thus this boundary is a tempered
Hilbert-valued distribution with the full ordered vector spectrum.
Taking both strong distributional boundaries in the reflected split,
first for factor tests and then for joint tests by continuity, gives
\[
 W(B^*A)=\langle\Psi(B),\Psi(A)\rangle
\]
for any two homogeneous words $A,B$ with their tests.  The concatenated
cross list is itself a list of homogeneous slots, even when the two
words have different total dimensions.

It remains to pass from homogeneous slots to arbitrary original
fields.  Write each such field in the fixed list as
$A_i=\sum_{\Delta\in S_i}A_{i,\Delta}$, with $S_i$ finite, as in
Assumption~\ref{ass:framework}.  Let $\mathbf A$ denote the complete
ordered list, including its $L$ slots, and let
$\mathbf A_{\boldsymbol\Delta}$ replace each old $A_i$ by
$A_{i,\Delta_i}$.  If there are no old slots, the indexing product
below is a singleton.  Multilinearity of the accessible analytic
products already gives the corresponding finite expansion in the
tube.  For an already constructed homogeneous list $\mathbf B$,
write $W_{\mathbf B}$ and $\Psi_{\mathbf B}$ for its scalar and vector
boundary maps.  For every joint Schwartz test $h$, in particular
every compact joint test, define the maps for $\mathbf A$ by
\begin{equation}\label{R:finite-field-expansion}
 \begin{split}
 W_{\mathbf A}(h)
   &=\sum_{\boldsymbol\Delta\in\prod_i S_i}
                 W_{\mathbf A_{\boldsymbol\Delta}}(h),\\
 \Psi_{\mathbf A}(h)
   &=\sum_{\boldsymbol\Delta\in\prod_i S_i}
                 \Psi_{\mathbf A_{\boldsymbol\Delta}}(h).
 \end{split}
\end{equation}
The same joint kernel $h$ occurs in every term.  This is a finite
expansion of field labels, not an assertion that $h$ has finite tensor
rank.  Each summand is an already constructed tempered distribution,
and every summand has the same cone for the chosen ordering.  Their
finite sum therefore has the required temperateness and support and
has the specified analytic product as its Fourier--Laplace continuation.
Adjunction and locality pass through the same finite expansion.
The exact scaling laws of the summands give similarity covariance with
dilations acting on the original field labels; a nonhomogeneous list
is not assigned one total dimension.  Orders and constants may be
chosen by taking a maximum of the finitely many seminorm orders and
a sum of the corresponding constants, without any uniform estimate
over dimensions or lengths.  Finite distributional derivatives in
any slot preserve temperateness and cone support and have the shifted
homogeneous covariance laws term by term.

The reflected identity also passes through this expansion with all
cross terms.  For arbitrary finite linear combinations $P,Q$ of
finite ordered words with joint tests, expand their original labels
finitely and apply the homogeneous split to
each pair of component words.  Here $P^*Q$ uses reversed adjunction
on each bra word and the tensor product of its conjugated, reversed
test with the ket test.  The result is
\begin{equation}\label{R:reflected-polynomials}
 W(P^*Q)=\langle\Psi(P),\Psi(Q)\rangle.
\end{equation}
In particular, the full joint Wightman positivity inequalities are
the Gram inequalities for these vectors in the original $\cH$.
This uses the cross-tuple splittings, not a claim of termwise
positivity for the individual summands.

The homogeneous induction, vector construction, and finite
identities \eqref{R:finite-field-expansion} and
\eqref{R:reflected-polynomials} complete the proof of
Theorem~\ref{R:mixed-products}.

\subsection{The joint domain and affiliation with the given net}

Let $\mathcal D$ be the span of the reconstructed vacuum products
with joint Schwartz test kernels, including compact joint kernels
and finite sums of factorized smearings.

The passage from positive vacuum distributions to an invariant
insertion domain is the Wightman reconstruction mechanism, in the
test-algebra formulation of Borchers~\cite{Wightman1956,Borchers1962}.
Here the vacuum vectors are already in the given $\cH$; after this
domain construction we must still prove affiliation in the fixed net.

For an old vacuum word represented by such a kernel, expand its field
labels as in
\eqref{R:finite-field-expansion} under the same joint kernel.
Old product multilinearity and the old-only identification in
Proposition~\ref{E:accessible-products} give equality of the old
word vector with this finite sum of homogeneous reconstructed word
vectors.  This exact label expansion makes no finite-rank assertion
about the joint kernel. Since $\mathcal D_{\mathrm{old}}$ is exactly the
algebraic span of these original vacuum-word vectors, it follows that
$\mathcal D_{\mathrm{old}}\subset\mathcal D$. In particular,
$\mathcal D$ is dense in the original $\cH$. This inclusion concerns the
specified vacuum-generated domain, not arbitrary additional vectors of
a separately chosen invariant domain.

Regard a test polynomial as a finite algebraic sum of finite ordered
words with the allowed kernels.  Multiplication concatenates the
words and tensors their kernels, and adjunction reverses the entries
and conjugates and reverses the kernel.  All resulting lengths remain
finite, with no fixed upper bound.  Left insertion defines joint
fields on $\mathcal D$.  To check that it is well-defined, let a whole
polynomial $A$ represent a zero vector.  Equation~\eqref{R:reflected-polynomials} gives
$W(A^*A)=0$.  Positivity and Cauchy--Schwarz applied to $A$ and
$B^*BA$, for any test polynomial $B$, give
\[
 |W(A^*B^*BA)|^2
 \leq W(A^*A)\,
       W\bigl((B^*BA)^*(B^*BA)\bigr)=0.
\]
Every term exists by the all-finite-length construction, and
$W(A^*B^*BA)=\|\Psi(BA)\|^2$.  Thus zero states form a left ideal.
This applies to the whole polynomial, even when its zero vector
results from cancellation between different scaling components.
Adjoint insertion gives
$A^\dagger(\bar f)|_{\mathcal D}\subset A(f)^*$, where the right
side is the Hilbert-space adjoint.  Its dense domain therefore makes
each smeared field closable.
For an old field inserted into a vector of
$\mathcal D_{\mathrm{old}}$, the same old-only identification applied to
the longer word, with the inserted test tensored with the original joint
kernel, proves agreement with the old insertion. For compact old smearings,
Section~\ref{sec:foundations} identified that insertion and its adjoint with
the restrictions of the original affiliated realizations. Thus the old fields and adjoints agree
with their original polynomial restrictions on
$\mathcal D_{\mathrm{old}}$, and
\eqref{R:leg-equation} becomes the operator identity
\begin{equation}\label{R:operator-leg}
 \Box L=\Theta\quad\hbox{on }\mathcal D.
\end{equation}
Translation covariance, similarity covariance, and locality are
operator identities on this common invariant domain.  No additional
Hilbert-space sector has been introduced.

This gives the common operator domain but has not yet supplied a
closed realization in a prescribed local algebra.  For that purpose
we prove exchange with every bounded element of that algebra's
commutant and then enlarge the domain by those commutant actions.
The distinction matters because weak field locality on
$\mathcal D$ alone does not establish affiliation.

We will use the same extension argument for the new scalar, the old
fields on the enlarged domain, and compact joint polynomials.  It is
convenient to prove it once for an adjoint pair.

\Needspace{6\baselineskip}
\begin{lemma}[Extension by commutant actions]\label{R:commutant-extension}
Let $\mathcal D$ be dense in $\cH$, and let $A,A^\sharp$ be operators
on $\mathcal D$.  Fix a localization set $K$ and the original
algebras $\cM(O)$ for open neighborhoods $O$ of $K$.  Suppose that
for every such $O$ and every $C\in\cM(O)'$,
\[
 \langle\phi,CA\psi\rangle
 =\langle A^\sharp\phi,C\psi\rangle,
 \qquad\phi,\psi\in\mathcal D,
\]
and that the corresponding identity also holds with $A,A^\sharp$
interchanged.  Then $A$ has a closed extension affiliated with every
$\cM(O)$, whose adjoint extends $A^\sharp$.  The reverse adjoint
pair has the same property.  If $A^\sharp=A$, the extension is
symmetric.
\end{lemma}

\begin{proof}
Let $\mathcal D_K$ be the span of $C_1\cdots C_m\psi$, with
$\psi\in\mathcal D$ and each $C_j$ in the commutant of an allowed
neighborhood algebra, including $m=0$.  Every finite collection of
these neighborhoods has a common smaller open neighborhood of $K$.
By isotony all its commutant coefficients belong to the commutant
of this smaller algebra.  Products and adjoints of the coefficients
therefore obey the assumed exchange identities.

For a finite sum with product coefficients $C_i$, prescribe
\[
 \widetilde A\sum_iC_i\psi_i=\sum_iC_iA\psi_i,
 \qquad
 \widetilde A^\sharp\sum_iC_i\psi_i=\sum_iC_iA^\sharp\psi_i.
\]
If the input sum is zero, pairing the proposed first image with
$\phi\in\mathcal D$ gives
$\langle A^\sharp\phi,\sum_i C_i\psi_i\rangle=0$.
Density proves well-definedness; the reverse exchange proves it
for the second prescription.  For two sums use exchange with
$D_j^*C_i$ in a common smaller algebra.  This gives
\[
 \widetilde A^\sharp\subset\widetilde A^*,\qquad
 \widetilde A\subset(\widetilde A^\sharp)^*.
\]
Both operators are thus closable.  Their common domain is invariant
under every unitary in every allowed commutant, and both operators
commute with that unitary and its inverse.  Their closed graphs retain
these invariances, proving affiliation simultaneously with all the
neighborhood algebras.  The adjoint inclusions survive closure.  In
the symmetric case the two prescriptions coincide, so their closure
is symmetric as well.
\end{proof}

\Needspace{6\baselineskip}
\begin{proposition}\label{R:affiliation}
For every real test $f$ with compact support in an open region $O$,
$L(f)$ has a closed symmetric extension affiliated with the original
algebra $\cM(O)$.  These realizations can be made
compatible for all neighborhoods of $\supp f$.
\end{proposition}

\begin{proof}
We first prove exchange with every bounded operator of the actual
commutant, not merely with opposite-region fields.  If
$C\in\cM(O)'$, then
\begin{equation}\label{R:commutant-exchange}
 \langle\phi,CL(f)\psi\rangle
       =\langle L(f)\phi,C\psi\rangle,
           \qquad \phi,\psi\in\mathcal D.
\end{equation}
First take $\phi,\psi$ to be finite sums of words with factorized
compactly supported tests, and by linearity treat one word on each
side.  Translate their slots into real mutually
spacelike patches, all spacelike from $\supp f$.  Subdivide $f$ into
local patches and choose its self-adjoint bounded approximants in
$O$.  The simultaneous approximating products converge there in
Hilbert norm.  Exchange holds at the bounded level because each
approximant to $L(f)$ commutes with $C$.  The norm limit proves
\eqref{R:commutant-exchange} on that real open placement set.

Continue separately in the translated polynomial slots on either
side.  A fixed $L(f)$ on the left of a polynomial leaves every suffix
spectral cone of that polynomial intact; on a bra the corresponding
opposite cones hold.  The bounded operator $C$ between the vectors
does not affect these two separate vector analyticity statements.
To see the common domain explicitly, use the positive vector tube
of \eqref{E:ordered-tube} for each polynomial before taking the bra
adjoint.  For a polynomial $\psi$ of length $n$, both $\psi(b)$ and
$L(f)\psi(b)$ are holomorphic on that tube.  In the latter product
let $q_1$ be its total momentum and
$q_2,\ldots,q_{n+1}$ its suffix momenta.  Fixing the first, $L$,
coordinate real removes its imaginary damping, but the strict
polynomial gaps still damp all $q_j$ with $j\geq2$.
The real smearing contributes $\widehat f(q_1-q_2)$, whose rapid
decay controls the remaining momentum since
$|q_1|\leq|q_1-q_2|+|q_2|$.  The tempered vector distribution
therefore has this smeared face, with all finite derivative bounds.
The same applies to $\phi$ and $L(f)\phi$.  After conjugating the
bra variables, the common scalar domain is the product of the
negative bra tube and the positive ket tube.  Boundedness of $C$
suffices for these pairings; neither commutation with translations
nor invariance of $\mathcal D$ under $C$ is used.
Their scalar pairings therefore obey boundary uniqueness, which
extends the identity to arbitrary placements and smearings.
For a compact joint kernel, first use a finite partition into product
patches and then approximate each piece by finite sums of factor
tests in the required finite smooth seminorms.  Theorem~\ref{R:mixed-products}
gives Hilbert-norm convergence both for the polynomial vector and
for the vector with the additional $L(f)$ slot.  Boundedness of $C$
therefore passes the identity to the joint kernel.  This test-topology
limit neither asserts finite tensor rank nor uses invariance of
$\mathcal D$ under $C$.  Compact joint tests are dense in the joint
Schwartz space, so the same Hilbert-norm continuity extends exchange
to all joint Schwartz kernels in $\mathcal D$.

Apply Lemma~\ref{R:commutant-extension} with
$A=A^\sharp=L(f)$ and $K=\supp f$.
The exchange identity holds for every neighborhood, so the lemma
gives one compatible closed symmetric affiliated extension.
Essential self-adjointness is neither asserted nor required.
\end{proof}

For a complex test $f=f_1+if_2$, linearity of the exchange identities
for the real tests gives the adjoint pair
$A=L(f)$, $A^\sharp=L(\bar f)$ on $\mathcal D$.
Lemma~\ref{R:commutant-extension} supplies a compatible closed
affiliated realization with the required adjoint.  This uses neither
self-adjointness nor strong commutation of the real and imaginary
parts.

We likewise need compatible realizations of the original entries on
the reconstructed domain before using polar tails on its vectors.
For an original field $A$ and compact test $f$, use the bounded
approximants of Proposition~\ref{E:bounded-approximation}, whose
vacuum and adjoint-vacuum convergence was established before the
mixed reconstruction.  Fix an open $O$ containing $\supp f$ and choose
these approximants $A_R(f)$ in $\cM(O)$.  Initially take the vectors
to be words with factorized compactly supported tests, as above.
At separated real placements the simultaneous
product limits give convergence of $A_R(f)\psi_R$ and
$A_R(f)^*\phi_R$.  Bounded commutation with $C\in\cM(O)'$ therefore
gives
\begin{equation}\label{R:old-commutant-exchange}
 \langle\phi,C A(f)\psi\rangle
   =\langle A^\dagger(\bar f)\phi,C\psi\rangle,
       \qquad \phi,\psi\in\mathcal D.
\end{equation}
Initially this holds on the separated placement set.  The same
separate-vector-tube argument, now with the first slot $A(f)$ or
$A^\dagger(\bar f)$ smeared over real coordinates, and the same
joint-test continuity
extend it to all the stated vectors.  Only the already proved mixed
vector spectra and the boundedness of $C$ enter this continuation.

Apply Lemma~\ref{R:commutant-extension} to the pair
$A(f),A^\dagger(\bar f)$ on $\mathcal D$, using
\eqref{R:old-commutant-exchange} and its adjoint version in every
neighborhood.  This gives compatible affiliated extensions of the
reconstructed restrictions, with adjoints extending the corresponding
insertions.  On $\mathcal D_{\mathrm{old}}$ these restrictions agree
with the original chosen realizations by the old-domain argument
above.  We do not require the newly reconstructed vectors to lie in
those previously chosen closed realizations, or identify their full
graphs with the extensions just constructed.

Only now use polar truncation for these realizations on polynomial
states.  If $\widehat A$ is the chosen closed realization, then
$\widehat A$ and $\widehat A^*$ agree on $\mathcal D$ with the
corresponding insertions, which preserve $\mathcal D$.  Induction gives
\[
 \mathcal D\subset
 \Dom(\widehat A^*\widehat A)^m
       \cap\Dom(\widehat A\widehat A^*)^m
 \qquad(m=1,2,\ldots),
\]
with the powers there equal to the alternating joint polynomial
words.  Their norms are finite by Theorem~\ref{R:mixed-products},
and spectral calculus gives the polar-tail estimates on every
polynomial vector.  For factorized polynomials, the natural-product
and polar-cutoff argument of Section~\ref{sec:foundations} applies to
these compatible realizations.  For compact joint kernels, take finite factor-test
approximations in the required test seminorms and pass both the vacuum
and adjoint-vacuum vectors by joint Hilbert-distribution continuity.
The closed wedge Tomita graph then passes its reality relation from
the factorized affiliated polynomial realizations to Hermitian joint
polynomial vectors.  These are the wedge reality properties needed
for the charge estimates of Section~\ref{sec:charges-cones}.

Finally, affiliation is only one part of the physical-membership
criterion in Assumption~\ref{ass:framework}.  The other part is also
present: $L=S$ is the explicit time-line extraction from actual
modular transforms of the physical virial field in arbitrarily small
diamonds.  Its separated products are the limits of those physical
bounded operations, and their joint continuation is uniquely fixed
by them.  Thus the field satisfies the stated intrinsic membership
criterion in the given net.  In a gauge formulation all these
operations take place in the chosen physical, gauge-invariant
algebras.

All construction requirements have now been checked: localization
in arbitrarily small original-net neighborhoods; joint field distributions
at every finite length on the common invariant domain
$\mathcal D\supset\mathcal D_{\mathrm{old}}$, with the required adjoints,
covariance, spectrum, and locality; and compatible affiliated smeared
realizations in the original net.
The notation $L$ henceforth denotes the resulting physical Hermitian
scalar of dimension two on $\mathcal D$.

\subsection{Improvement and the local Ward algebra}

Define on $\mathcal D$
\begin{equation}\label{R:improvement}
 t_{\mu\nu}=T_{\mu\nu}
       +\frac13(\partial_\mu\partial_\nu
                        -\eta_{\mu\nu}\Box)L.
\end{equation}
It is a physical symmetric conserved tensor, and
\eqref{R:operator-leg} gives $t^\mu{}_{\mu}=0$.
For an affine Poincar\'e Killing vector $X$, its charge difference
from $T$ is the divergence of the antisymmetric superpotential
\begin{equation}\label{R:superpotential}
 \frac13\partial^\lambda
 \left(X_\lambda\partial_\mu L-X_\mu\partial_\lambda L
                              +(\partial_\lambda X_\mu)L\right).
\end{equation}
Its compact commutator flux vanishes, so translation and Lorentz
charges agree.  If the initial stress tensor required the weight replacement of
Section~\ref{sec:foundations}, also include the local difference
between that selected tensor and the initial one.  That difference
is symmetric and conserved and has zero affine charges by the
weight-replacement argument.  This proves the improvement criterion.  We now check the special-current Ward algebra on the
actual reconstructed domain.

We first construct the local Ward action on the represented polynomial
algebra generated by the original fields and $L$; this is distinct from
the vector space $\mathcal F$ of point-field labels.
Proposition~\ref{R:compatible-families} will extend the action to every
further jointly admitted finite family.  For an open $O$, let
$\mathfrak P_c(O)$ consist of
finite sums of finite ordered products of original fields, $L$, and their
covariant derivatives, smeared by compact smooth joint kernels supported
in $O^n$, and identified when their restrictions to $\mathcal D$ agree.
The joint-domain construction makes concatenation and reversed adjunction
well-defined on this quotient and leaves $\mathcal D$ invariant.  Write
$P^\dagger$ for the reversed field-adjoint polynomial, to distinguish it
from the Hilbert adjoint.  Every finite length and every admitted compact
joint kernel are included; the kernel need not have finite tensor rank.

These compact polynomial insertions have compatible affiliated
realizations in the original net.  For a factorized word, take the
natural-domain product of the single-field realizations constructed
above, and take natural sums for finite polynomials.  Their domains
contain $\mathcal D$, and their adjoints contain the reversed adjoint
polynomials there, so they are closable.  Every commutant unitary and
its inverse preserve their natural graphs; their closures are therefore
affiliated.  This uses no commutation among the unbounded factors.
In particular, $t(f)$ is the finite-sum case of this argument.

For a compact joint kernel, partition into product patches inside $O^n$
and approximate in the smooth test topology by finite sums of tensor
tests.  If $P_j$ are the resulting factorized polynomials, then
Theorem~\ref{R:mixed-products}, with any fixed $\mathcal D$ word
appended, gives
\[
 P_j\psi\longrightarrow P\psi,\qquad
 P_j^\dagger\phi\longrightarrow P^\dagger\phi
 \quad(\phi,\psi\in\mathcal D)
\]
in Hilbert norm.  Affiliation of the factorized closures gives exchange
with commutant unitaries, hence with every $C\in\cM(O)'$ by linear
combinations of unitaries.  Boundedness of $C$ passes it to the limit:
\[
 \langle\phi,CP\psi\rangle
       =\langle P^\dagger\phi,C\psi\rangle,
       \qquad \phi,\psi\in\mathcal D.
\]
The reversed adjoint polynomial obeys the reverse exchange identity.
Lemma~\ref{R:commutant-extension}, with $A=P$ and
$A^\sharp=P^\dagger$, therefore supplies compatible affiliated
extensions for every neighborhood of the compact localization.
This is an extension with the required adjoint, not an assertion
about the minimal closure or a limit of the factorized closures.
Its polar truncations are bounded original-net operations converging on
$\mathcal D$, with adjoint convergence there.  All the commutants here
are those of the original net.

For a real conformal Killing field $X$, set
$j_{X,\mu}=t_{\mu\nu}X^\nu$ and
$\sigma_X=(\partial\cdot X)/4$.  Conservation and tracelessness give
$\partial^\mu j_{X,\mu}=0$.  Define
\begin{equation}\label{R:Ward-definition}
 \delta_X A=i[q_X,A]
\end{equation}
as a compact commutator flux enclosing the insertion.  This is an
operator insertion on $\mathcal D$, without assuming the existence
of a global unbounded charge on every state.  More explicitly, smear
the two fields with a compact kernel in their relative coordinate
and the target test in their merged coordinate.  Caps carry the same
unit-flux normalization and orientation.  In relative coordinates the
commutator current is divergence-free and supported in causal positions.
The difference of two normalized enclosing caps is a compact gradient
there, up to spacelike outer cutoff terms.  Conservation and locality
therefore allow exact deformation or shrinking while keeping the target
enclosed.  The resulting expression is a finite sum of joint smeared
products; combining the two momenta preserves the ordered cones for
the remaining slots.  Polynomial coefficients of $X$ give finite
derivatives of test kernels and do not spoil these properties.

For a joint target compactly localized in $O$, partition into product
patches and choose the relative caps within their localization margins.
A flux about a word is the finite sum of the fluxes about its slots, so
these representatives belong to $\mathfrak P_c(O)$.  To check the quotient
by zero operators, choose one absolute normalized compact cap enclosing
the finite union of target projections for the polynomial representatives.
It may lie outside $O$.  Deformation identifies every slotwise cap with
this cap.  The algebraic commutator rule identifies the sum with
$i[Q_X,P]$ on factor tests, and the same joint Hilbert continuity
extends the identity to the joint kernel.  Since $Q_X$ and $P$ preserve $\mathcal D$,
this vanishes if $P|_{\mathcal D}=0$.  Thus the local derivations are
well-defined on $\mathfrak P_c(O)$ and preserve it; their iterates are
finite joint polynomials.  The preceding exchange argument supplies
their compatible original-net affiliated extensions.

The complete vacuum flux vanishes on compact joint products by
Lemma~\ref{C:low-energy} and its subsequent dyadic charge estimate.
That proof applies here because the new compact polynomial vectors
have the wedge reality properties just proved, while the current
vacuum wave is the same purely timelike, translation-smooth wave
used there.  Its mass-band estimate is uniform; the summable
low-energy and high-energy dyadic bounds remain unchanged.
Consequently
\begin{equation}\label{R:vacuum-Ward}
 \sum_{i=1}^n
 W(A_1,\ldots,\delta_XA_i,\ldots,A_n)=0.
\end{equation}
The derivation preserves adjunction.  Its translation, Lorentz, and
dilation parts on the original fields are the physical actions:
\eqref{R:superpotential} handles the first two, and
Proposition~\ref{C:Q-zero} removes the possible conserved virial
charge in the third.  Identities first proved on an old vacuum leg
extend to $\mathcal D$ by locality, old polynomial density, and
ordered boundary uniqueness.

We need the transformation of the current itself to identify the
Lie algebra.  The two required primary identities are
\begin{align}
 \delta_X L&=X^\lambda\partial_\lambda L+2\sigma_XL,
                         \label{R:primary-L}\\
 \delta_X t_{\mu\nu}
   &=X^\lambda\partial_\lambda t_{\mu\nu}
      +(\partial_\mu X^\lambda)t_{\lambda\nu}
      +(\partial_\nu X^\lambda)t_{\mu\lambda}
      +2\sigma_Xt_{\mu\nu}. \label{R:primary-t}
\end{align}
We explain their return from the previously established pair
representation.  All short analytic products now agree with the
joint tempered ones: they agree on separated vacuum-wave products,
then everywhere by continuation.  Thus the Wick sphere/cap
comparison of Proposition~\ref{P:flux-wards} applies directly to
physical joint pair boundaries.  For a translation and target $T$,
the Lorentzian cap difference is the physical translation Ward
identity.  Its sphere version has the common nonzero normalization
fixed there.  The translation sphere identity for $L$ was already
proved, and differentiating it at a fixed enclosing surface gives
the same identity for the derivative terms in
\eqref{R:improvement}.  It therefore holds with target $t$.

Conjugate these translation identities in the positive pair
representation of Proposition~\ref{P:pair-representation}, whose
primary targets include $L$ and $t$.  Transporting the surrounding
sphere and its holomorphic Killing coefficient gives the sphere
identities for every conformal Killing field.  Return to Lorentzian
commutators with the same averaged cylinders as in
Proposition~\ref{P:flux-wards}.  The time caps converge by the joint
tempered boundary theorem; the sides have a fixed spacelike margin
and vanish with their thickness.  The Wick pullback of $X$ is
holomorphic on these compact surfaces, so the identical argument
works for its degree-zero, degree-one, and degree-two coefficients.
The normalization is the translation normalization in every case.
This proves \eqref{R:primary-L}--\eqref{R:primary-t} on the vacuum.
Their discrepancies are relatively local insertions.  Move such an
insertion to an extreme position on a spacelike open set; old
polynomial density and ordered uniqueness then show that its vacuum
vanishing implies vanishing on $\mathcal D$.  The two displayed
identities are therefore operator insertion identities on that
domain.

\Needspace{6\baselineskip}
\begin{proposition}\label{R:Ward-Lie}
The compact-flux Ward derivations realize the conformal Lie algebra
with the active-action convention
\begin{equation}\label{R:Lie-sign}
 [\delta_X,\delta_Y]=-\delta_{[X,Y]}.
\end{equation}
They obey \eqref{R:vacuum-Ward} and extend the physical affine Ward
derivations on the original field algebra.
\end{proposition}

\begin{proof}
Coordinates multiplying fields are held fixed when taking a Ward
variation.  Using \eqref{R:primary-t}, the conformal Killing equation,
and conservation, direct differentiation gives
\begin{equation}\label{R:current-variation}
 \delta_X j_{Y,\mu}
  =\partial^\lambda(X_\lambda j_{Y,\mu}
                              -X_\mu j_{Y,\lambda})
       -j_{[X,Y],\mu}.
\end{equation}
For example, expanding the divergence gives
$4\sigma_X j_{Y,\mu}+X\cdot\partial j_{Y,\mu}
-(\partial^\lambda X_\mu)j_{Y,\lambda}$; subtracting
$j_{[X,Y]}$ cancels the derivative of $Y$, and
$\partial_\lambda X_\mu+\partial_\mu X_\lambda
=2\sigma_X\eta_{\lambda\mu}$ yields exactly the two index terms
and the remaining $2\sigma_X$ in \eqref{R:primary-t}.

Choose one compact flux $Q_Y(F)$ implementing $\delta_Y$ on an
insertion $A$ and its local $X$ variation.  Use a larger compact
flux for $\delta_X$ on $Q_Y(F)$ and these insertions.  The derivation
rule on the joint domain gives
\[
 [\delta_X,\delta_Y]A=i[\delta_X Q_Y(F),A].
\]
The first term on the right of \eqref{R:current-variation} is an
antisymmetric divergence.  Its commutator flux about the compact
insertion is zero: integration by parts leaves only the outer
boundary, where locality permits the cutoff to be removed.
The second term therefore gives
$-i[Q_{[X,Y]}(F),A]=-\delta_{[X,Y]}A$.
All products and iterated insertions in this calculation are defined
by compact joint kernels on $\mathcal D$.  Equation
\eqref{R:vacuum-Ward} and the affine identification were proved
above.
\end{proof}

\subsection{Centered fields and compatible families}

A Ward variation with polynomial coordinate coefficients is an insertion,
but need not itself be a translation-covariant point-field label.  We now
remove that coordinate part and verify the intrinsic membership criterion
for each finite family needed in a finite Ward calculation.  Write
\[
 \begin{gathered}
 p_c(z)=c,\qquad d(z)=z,\qquad
 m_{b,c}(z)=(b\cdot z)c-(c\cdot z)b,\\
 k_b(z)=2(b\cdot z)z-z^2b,\qquad b,c\in\R^{1,3}.
 \end{gathered}
\]
Fix a homogeneous finite Lorentz multiplet $A$ of dimension $\Delta$
among the original fields, $L$, and their covariant derivatives.
Let $N_A$ be the dimension of its component space.  Define the centered insertion
\[
 B_{b,A}(x)=\delta_{k_b^{[x]}}A(x),\qquad
 k_b^{[x]}(z)=k_b(z-x).
\]
For a target test $f$, this denotes $i$ times the difference of joint products
with relative kernel $f(x)h^\mu(r)k_b^\nu(r)$, $r=y-x$, where $h$ is a
normalized enclosing cap.  It involves no pointwise product or
unsmeared surface restriction.  The cap deformation just proved makes
the definition independent of its admissible normalized choice.

Define its iterates on the same polynomial domain by
\[
 \begin{aligned}
 B_A^{(0)}&=A,\\
 B^{(r)}_{b_r\ldots b_1,A}(x)
   &=\delta_{k_{b_r}^{[x]}}
       B^{(r-1)}_{b_{r-1}\ldots b_1,A}(x),\qquad r\ge1.
 \end{aligned}
\]
Every fixed expression is a finite compact joint polynomial because
\(\delta_X\) preserves the local polynomial algebra.

\Needspace{6\baselineskip}
\begin{proposition}[Finite families of centered fields]
\label{R:centered-families}
Let $A$ range over a fixed finite collection of homogeneous original
fields, $L$, and their covariant derivatives.  For every fixed finite
collection of depths, the corresponding centered fields $B_A^{(r)}$
have all mutual and original-field joint products on $\mathcal D$,
with the required adjoints, locality, covariance, ordered spectra,
and compatible original-net affiliated realizations.  If $A$ has
real dimension $\Delta$ and $N_A$ components, depth $r$ has dimension
$\Delta-r$ and at most $4^rN_A$ components before field identities.
These finite families satisfy the intrinsic membership criterion of
Assumption~\ref{ass:framework}.  Every finite Ward word on the
original fields and $L$ is a finite sum of polynomial-coordinate
multiples of derivatives of such centered fields.
\end{proposition}

\begin{proof}

For a similarity $g(z)=\lambda\Lambda z+a$, $\lambda>0$, write $U(g)$
for the already given Poincar\'e/dilation implementer.  Changing variables
in the smooth current test for the dimension-four $t$ sends
$X$ to $g_*X=\lambda\Lambda X\circ g^{-1}$ and an absolute normalized
cap test $h_0$ to $(h_0)_g=\lambda^{-1}\Lambda h_0\circ g^{-1}$,
in the tensor convention of their contraction.  For a relative cap,
the translation cancels in this formula.  Since
$g_*k_b^{[x]}=\lambda^{-1}k_{\Lambda b}^{[gx]}$, cap independence gives
\[
 U(g)B_{b,A}(x)U(g)^{-1}
   =\lambda^{\Delta-1}\rho_A(\Lambda)
                       B_{\Lambda b,A}(gx),
\]
where $\rho_A$ is the Lorentz component matrix of $A$.  Thus $B_{b,A}$
is a possibly zero translation-covariant homogeneous insertion of real
grade $\Delta-1$, in a finite Lorentz family with one additional vector index.
Only the existing similarity covariance and the proved cap deformation
have been used.

Let $s_{b,c}$ denote the constant component matrix in the affine action
\[
 \delta_{p_c}A=c\cdot\partial A,\qquad
 \delta_dA=(x\cdot\partial+\Delta)A,\qquad
 \delta_{m_{b,c}}A=(m_{b,c}(x)\cdot\partial+s_{b,c})A.
\]
Expanding
$k_b(z-x)=k_b(z)-2(b\cdot x)d(z)-2m_{b,x}(z)+p_{k_b(x)}(z)$
in the centered flux gives
\begin{equation}\label{R:centered-Ward}
 \delta_{k_b}A(x)
 =\bigl(k_b(x)\cdot\partial+2\Delta(b\cdot x)+2s_{b,x}\bigr)A(x)
       +B_{b,A}(x).
\end{equation}
The polynomial-coordinate term is an insertion, in general not a
translation-covariant point-field label; its coefficients multiply
the target test after smearing.  For example,
\eqref{R:primary-L} gives $B_{b,L}=0$.

The local affine Ward actions on $B$ also follow from the proved Lie
relation.  Indeed
$[p_c,k_b]=2(b\cdot c)d+2m_{b,c}$ and $[d,k_b]=k_b$.
At a fixed enclosing cap target differentiation commutes with the
commutator, so \eqref{R:Lie-sign} gives
\[
 \begin{aligned}
 (\partial_c-\delta_{p_c})\delta_{k_b}A
    &=2(b\cdot c)\delta_dA+2\delta_{m_{b,c}}A,\\
 \delta_d\delta_{k_b}A
    &=(x\cdot\partial+\Delta-1)\delta_{k_b}A .
 \end{aligned}
\]
The polynomial term in \eqref{R:centered-Ward} has the same two
defects.  Subtraction gives
$\delta_{p_c}B=c\cdot\partial B$ and
$\delta_dB=(x\cdot\partial+\Delta-1)B$.
The Lorentz bracket $[m_M,k_b]=-k_{Mb}$ for $m_M(z)=Mz$ similarly
gives the tensor-product component action.  Thus these affine
identifications are consequences, not an assumed conformal closure.

For any fixed finite collection of centered insertions, substituting
their relative kernels into a joint product expands it into finitely
many of the already constructed products.  This is a continuous map
of compact smooth or Schwartz test spaces: the relative supports are
compact and the coefficients are polynomial.  It supplies all joint
products at every finite length on $\mathcal D$, including adjoints
and all mixed orderings.  Merging each block retains the ordered suffix
cones.  For real $b$ and $h$, reversing the two products gives
$B_{b,A}^\dagger=B_{b,A^\dagger}$.  Positivity holds in the original
Hilbert space, and locality
follows by shrinking the relative supports on spacelike product patches;
a compact test on the spacelike configuration set is first partitioned
into such patches.  The preceding exchange argument supplies compatible
original-net affiliated realizations and their adjoints for these finite
homogeneous families.

At each depth, choose a shrinking representative of the previous depth
inside an inner neighborhood, and its enclosing cap inside the prescribed
target neighborhood.  The previous depth's locality makes its current
commutator causally supported, so conservation gives the same normalized-cap
deformation and shrinking property at the next depth.  Every nested
expression is still a finite joint polynomial in the original fields
and $L$.  The joint continuity and full-commutant construction above
therefore give all mutual products and compatible affiliation for any
fixed finite collection of depths.  The covariance calculation gives
grade $\Delta-r$ and at most $4^rN_A$ components before field identities.
Equation~\eqref{R:centered-Ward}, the affine actions just
proved, and finite Leibniz rules express every fixed finite Ward word
as a finite sum of polynomial-coordinate multiples of finite derivatives
of these families.  For an original
$A=\sum_{\Delta\in S_A}A_\Delta$, do this separately for the finitely
many components, using the same target joint kernel in every component
term and assigning no single weight to their sum.  Thus each finite Ward
calculation uses homogeneous insertions in a jointly constructed finite
family on $\mathcal D$.  This does not identify them with original labels
in $\mathcal F$, assert global closure under centered descendants, or
construct a conformal action on the weakly closed bounded net.
\end{proof}

The same currents also act on other fields already admitted by the
intrinsic criterion.  This requires a joint reconstruction with $L$;
separate admission of two families is not by itself a statement about
their mutual products.

\Needspace{6\baselineskip}
\begin{proposition}[Extension to compatible physical families]
\label{R:compatible-families}
Let $\mathcal G$ be a finite homogeneous finite-component family
jointly admitted with the original fields under
Assumption~\ref{ass:framework}.  The scalar $L$ has all joint products
with this family and the original fields, with compatible adjoints and
original-net affiliated realizations on a common invariant domain
$\mathcal D_{\mathcal G}\supset\mathcal D$ in the original Hilbert
space.  On the resulting represented polynomial algebra, the same
currents $j_X$ give unbroken local conformal Ward derivations,
extending the physical Poincar\'e and dilation actions on
$\mathcal G$ as well as on the original fields.  Every fixed finite
family of their centered Ward descendants satisfies the same intrinsic
membership criterion.  These constructions agree when one compatible
family is contained in another.
\end{proposition}

\begin{proof}
Regard the original fields and $\mathcal G$ as the spectator fields
in the boundary reconstruction of this section.  Their joint products,
common invariant realization, adjoints, and original-net affiliations
are given by the admission criterion.  Each fixed homogeneous list
has tempered vacuum distributions by the scaling argument of
Section~\ref{sec:foundations}; finite sums treat the original labels.
Their affiliated polar truncations obey
Proposition~\ref{E:bounded-approximation}, since the required moments
are alternating words in these already given joint products.  Thus
Proposition~\ref{E:accessible-products} supplies their analytic products
with the same accessible $L$ and the same reflected splittings.
No current Ward identity on $\mathcal G$ is needed for this step.

The pure $L$ construction is unchanged.  In the simultaneous mixed
induction, replacing $L$ by $\Theta$ lowers the number of $L$ entries
and adds an original homogeneous spectator.  The merged-face estimate,
causal ambiguity, Gaussian coefficient recovery, and scaling argument
therefore apply with $\mathcal G$ among the fixed spectator labels.
They prove Theorem~\ref{R:mixed-products} for every such enlarged
list.  Its positivity and full-commutant exchange then give the
invariant joint domain $\mathcal D_{\mathcal G}$ and the compatible
affiliated realizations exactly as above.  Products involving only
the original fields and $L$ agree with the previous ones: their
analytic products are the same approximation-independent limits,
and their tempered boundaries are unique.  Consequently
$\mathcal D\subset\mathcal D_{\mathcal G}$ and all insertion actions
agree there.  Old-and-$\mathcal G$ products likewise retain their
given vacuum-generated restrictions, without requiring agreement
of chosen closed realizations on an independently larger domain.
The same reasoning proves agreement for nested
compatible families.  If derivatives of the added fields are included
as labels, their dimensions lie in the finite union of progressions
$\Delta+\mathbb Z_{\ge0}$; each fixed grade receives only finitely
many derivative components.  This preserves finite scaling support
and the stipulated locally finite, bounded-below grade convention.

Compact commutator fluxes now define $\delta_X$ on the enlarged
polynomial algebra.  Conservation, locality, and joint continuity
give cap deformation, the derivation rule, and preservation of the
quotient by zero operators just as before.  Its compact polynomial
vectors have wedge reality by their affiliated realizations and the
joint graph limits.  Lemma~\ref{C:low-energy} and its summable dyadic
estimate therefore give vanishing complete vacuum flux and the Ward
identity \eqref{R:vacuum-Ward} on the enlarged domain.

We check the physical affine actions on the added fields, rather
than assume them.  Let $X$ be a translation, Lorentz, or dilation
generator, and let $G_X$ be its given self-adjoint generator, with
the convention that the physical variation is $i[G_X,\cdot]$.
For a compact smearing $F$ of a member of $\mathcal G$ and an old
compact polynomial $C$, the vacuum Ward identity gives
\[
 \langle C\Omega,\delta_XF\,\Omega\rangle
 =-\langle(\delta_XC)\Omega,F\Omega\rangle
 =i\langle G_XC\Omega,F\Omega\rangle
 =\langle C\Omega,iG_XF\Omega\rangle.
\]
The middle identity uses the already proved affine action on old
polynomials.  Smooth compact smearing and the given similarity
covariance place $F\Omega$ and $C\Omega$ in the generator domains;
the vacuum is invariant.  Density of old polynomial vectors thus
identifies the vacuum action with the physical affine variation of
$F$, including its component and weight terms.

The discrepancy between these two variations is relatively local
and closable on $\mathcal D_{\mathcal G}$: it is a finite compact
joint insertion with a densely defined reversed adjoint.  It
annihilates the vacuum.  Choose an old double cone spacelike to its
compact localization.  Locality makes it annihilate every old
polynomial vacuum vector from that cone, a dense set by
Reeh--Schlieder.  The closure of its graph therefore contains
$\cH\times\{0\}$, and being an operator graph it can contain no
nonzero output.  The discrepancy vanishes on the whole joint domain.
This proves the required affine identifications without assuming
new current Ward identities.

The primary identities for $L$ and $t$ likewise persist on the
enlarged domain: their relatively local discrepancies vanish on the
original dense domain, and the same adjoint and closed-graph argument
applies.  The current-variation computation and compact nested flux
proof of Proposition~\ref{R:Ward-Lie} consequently give the conformal
Lie relation on every enlarged polynomial insertion.  Finally the
centered construction of Proposition~\ref{R:centered-families} uses
only these Ward identities, similarity covariance, cap shrinking,
and finite joint products.  Repeating it on
$\mathcal D_{\mathcal G}$ proves its conclusion for every fixed
finite family of added targets and finite descendant depths.
\end{proof}

Thus the improvement \eqref{R:improvement} lies in the original
physical theory, has the original translation and Lorentz charges,
and is traceless.  Its special-conformal currents are physical
conserved currents with unbroken local Ward identities.  The
conclusion concerns flat spacetime and does not impose a
curved-background Weyl identity.
